\documentclass[aps,pra,11pt,onecolumn,nofootinbib,tightenlines]{revtex4-2}

\usepackage{mathrsfs}
\usepackage{amsmath}
\usepackage{amsfonts}
\usepackage{amssymb}
\usepackage{amsthm}
\usepackage{bm}
\usepackage[shortlabels]{enumitem}
\usepackage[colorlinks=true,urlcolor=blue,citecolor=blue,linkcolor=blue]{hyperref}
\usepackage{xurl}
\usepackage{natbib}

\setlength{\jot}{7pt}
\allowdisplaybreaks

\newtheorem{theorem}{Theorem}[section]
\newtheorem{lemma}[theorem]{Lemma}
\newtheorem{proposition}[theorem]{Proposition}
\theoremstyle{definition}
\newtheorem{remark}[theorem]{Remark}

\newcommand{\mc}[1]{\mathcal{#1}}
\newcommand{\mb}[1]{\mathbb{#1}}
\newcommand{\R}{\mathbb{R}}
\newcommand{\Y}{\mathcal{Y}}
\def\d{\mathrm{d}}
\renewcommand{\thesection}{\arabic{section}}
\makeatletter
\renewcommand\@pnumwidth{2.5em}
\makeatother

\begin{document}

\title{Sufficient and Necessary Conditions for the Parameters of Conditional Entropies}

\begin{abstract}
This standalone note gives the sufficient and necessary parameter conditions for the
conditional entropies introduced in Ref.~\cite{rubboli2026complete}.  It is a concise,
human-readable replacement for Sections~4 and~5 of that paper.  The statements cover
finite and general measures, both tropical endpoints, and derivations.  In particular,
no global finiteness assumption is imposed on the measure.  The proofs retain the main
ideas---finite-support factorization, common discretization, and localized
high-dimensional counterexamples---while omitting the lengthy analytic calculations.
\end{abstract}

\author{Roberto Rubboli}
\noaffiliation
\author{Erkka Haapasalo}
\noaffiliation
\author{Marco Tomamichel}
\noaffiliation

\maketitle

\tableofcontents

\section{Sufficient conditions for the parameters}
\label{sec:sufficient-conditions-measure}

In this section, we establish sufficient conditions for the parameters under which the
homomorphisms and derivations (and, consequently, the associated conditional entropies)
are monotone under conditional majorization.  In the following section, we show that
these conditions are also necessary, without imposing a separate regularity assumption
on the measure.

In Ref.~\cite{rubboli2026complete}, the general forms of the homomorphisms and
derivations depend on a finite measure: in the temperate case one may write it as
$\mu=t\tau$, with $\tau$ a probability measure and its sign encoded by $t$.
However, not all such choices are valid, as some lead to violations of monotonicity
under conditionally mixing channels.
While monotonicity under mixing channels acting on the first register always holds,
monotonicity under channels acting on the conditioning register fails for some choices
of the parameters.  We therefore examine the convexity properties of the homomorphism
as a function of the joint probability matrix.  Proposition~A.1 and Lemma~A.3 of
Ref.~\cite{rubboli2026complete} identify these convexity properties with the required
monotonicity and lift the columnwise result to the original semiring and channel
preorder of Eqs.~(43) and~(49) therein.

Let us first identify the function for which we need to establish convexity or
concavity properties.  In the temperate case, Eq.~(78) and Proposition~3.6 of
Ref.~\cite{rubboli2026complete} give
\begin{align}
\label{eq:TemperateForm}
\Phi(P_{XY})&=\Phi_{t,\tau}(P_{XY})
=\sum_{y\in\Y}P_Y(y)
  \exp\!\left(t\int_{[0,+\infty]}H_\alpha(P_{X|Y=y})\,
  \d\tau(\alpha)\right),
\end{align}
where $\tau$ is a probability measure and
$t\in\R\setminus\{0\}$.  Thus it suffices to study, for
$\vec{x}\in\R^d_{\geq0}\setminus\{0\}$,
\begin{equation}
\label{eq:conditionfunct}
\vec{x}\longmapsto
\|\vec{x}\|_1\exp\!\left(
t\int_{[0,+\infty]}H_\alpha(\hat{x})\,\d\tau(\alpha)
\right),
\qquad
\hat{x}=\frac{\vec{x}}{\|\vec{x}\|_1}.
\end{equation}
The R\'enyi family and its endpoint values $H_0$, $H_1$, and $H_\infty$ are those
of Eq.~(2) and the sentence following it in Ref.~\cite{rubboli2026complete}.
We begin with finite support, pass to a general
measure by a common discretization, and finally take the tropical and derivation
limits.

\subsection{Sufficient conditions for the discrete case}

The discrete case exploits two basic facts.  The first is Minkowski's inequality
\cite[Theorem~9]{bullen2013handbook}.

\begin{lemma}
\label{lem:Minkowski}
Let $\alpha>0$.  The function, defined on
$\R^d_{>0}$,
\begin{equation}
\vec{x}\longmapsto\left(\sum_{i=1}^d x_i^\alpha\right)^{1/\alpha},
\end{equation}
\begin{enumerate}[itemsep=1pt]
\item is convex if $\alpha>1$;
\item is concave if $0<\alpha<1$.
\end{enumerate}
\end{lemma}

The displayed power mean is not defined at $\alpha=0$.  Whenever the entropy
parameter equals $0$, the proof instead uses the endpoint formula for $H_0$.

The second fact is the standard convexity criterion for multivariate monomials
\cite[Lemma~8]{mu21_renyi}; see also Ref.~\cite[Proposition~13]{farooq2024matrix}.

\begin{lemma}
\label{lem:Matrix-majorization}
Let $N\in\mb{N}$ and let
$\vec{\beta}=(\beta_1,\ldots,\beta_N)\in\R^N$ satisfy
$\sum_{\ell=1}^N\beta_\ell=1$.  The function, defined on
$\R^N_{>0}$,
\begin{equation}
\vec{y}\longmapsto\prod_{\ell=1}^N y_\ell^{\beta_\ell},
\end{equation}
\begin{enumerate}[itemsep=1pt]
\item is concave if and only if $\beta_\ell\geq0$ for every $\ell$;
\item is convex if and only if there is a unique index $k$ such that
$\beta_k>0$ and $\beta_\ell\leq0$ for every $\ell\neq k$.
\end{enumerate}
\end{lemma}

Combining these facts gives the finite-support result.

\begin{lemma}
\label{lem:discrete-measure}
Let $\tau$ be a discrete probability measure supported on finitely many points
$\alpha_1<\cdots<\alpha_N$ of $[0,+\infty]$, and let
$t\in\R\setminus\{0\}$.  For every $d\in\mb{N}$, the function
\begin{equation}
\vec{x}\longmapsto\|\vec{x}\|_1
\exp\!\left(t\int_{[0,+\infty]}H_\alpha(\hat{x})\,
\d\tau(\alpha)\right)
\end{equation}
on $\R^d_{\geq0}\setminus\{0\}$ is
\begin{enumerate}
\item concave if $t>0$, $\alpha_1,\ldots,\alpha_N<1$, and
\begin{equation}
\sum_{\ell=1}^N\frac{\alpha_\ell}{1-\alpha_\ell}
\tau(\{\alpha_\ell\})\leq\frac1t;
\end{equation}
\item convex if $t<0$ and either
  \begin{enumerate}[label=(\alph*)]
  \item $\alpha_1,\ldots,\alpha_N\leq1$, or
  \item $\alpha_1,\ldots,\alpha_{N-1}<1$, $\alpha_N>1$, and
  \begin{equation}
  \sum_{\ell=1}^N\frac{\alpha_\ell}{1-\alpha_\ell}
  \tau(\{\alpha_\ell\})\leq\frac1t,
  \end{equation}
  where $\alpha/(1-\alpha)$ has value $-1$ at $\alpha=+\infty$.
  \end{enumerate}
\end{enumerate}
\end{lemma}

\begin{proof}[Proof idea]
First assume that every atom satisfies
$0<\alpha_\ell<+\infty$ and $\alpha_\ell\neq1$, and set
\begin{equation}
\beta_\ell=t\frac{\alpha_\ell}{1-\alpha_\ell}
\tau(\{\alpha_\ell\}),
\qquad
\bar\beta=1-\sum_{\ell=1}^N\beta_\ell.
\end{equation}
The function factors as
\begin{equation}
\left(\sum_i x_i\right)^{\bar\beta}
\prod_{\ell=1}^N
\left(\sum_i x_i^{\alpha_\ell}\right)^{\beta_\ell/\alpha_\ell}.
\end{equation}
Lemma~\ref{lem:Minkowski} supplies the convexity or concavity of each power mean,
and Lemma~\ref{lem:Matrix-majorization} supplies the joint sign criterion for the
outer monomial.  When $t<0$ and all support lies in $[0,1]$, no moment bound is
needed.  Atoms at $0$, $1$, and $+\infty$ are handled by the endpoint formulas
and endpoint-preserving perturbations or pointwise limits; the branch containing
$\alpha=1$ occurs only in the all-lower negative case.
\end{proof}

\subsection{Sufficient conditions for the temperate homomorphisms - general measure}

We now extend the result to a general probability measure $\tau$.  The signed
integral
\begin{equation}
\label{eq:IntegralCond}
\int_{[0,1)}\frac{\alpha}{1-\alpha}\,\d\tau(\alpha)
+\int_{(1,+\infty]}\frac{\alpha}{1-\alpha}\,\d\tau(\alpha)
\leq\frac1t
\end{equation}
is used only when the two one-sided integrals are finite; at $+\infty$ the
integrand is $-1$.  In particular, Eq.~\eqref{eq:IntegralCond} never hides an
undefined cancellation at $\alpha=1$.

\begin{proposition}
\label{prop:sufficient-conditions}
Let $\tau$ be a probability measure on $[0,+\infty]$ and let
$t\in\R\setminus\{0\}$.  The function in Eq.~\eqref{eq:conditionfunct} is
\begin{enumerate}
\item concave if $t>0$,
$\operatorname{supp}(\tau)\subseteq[0,1]$, $\tau(\{1\})=0$, and
\begin{equation}
\int_{[0,1)}\frac{\alpha}{1-\alpha}\,\d\tau(\alpha)\leq\frac1t;
\end{equation}
\item convex if $t<0$ and either
  \begin{enumerate}[label=(\alph*)]
  \item $\operatorname{supp}(\tau)\subseteq[0,1]$, with no moment-finiteness
  condition, or
  \item there is $\alpha_*\in(1,+\infty]$ such that
  $\operatorname{supp}(\tau)\subseteq[0,1]\cup\{\alpha_*\}$,
  $\tau(\{1\})=0$, the two one-sided integrals in
  Eq.~\eqref{eq:IntegralCond} are finite, and Eq.~\eqref{eq:IntegralCond} holds.
  \end{enumerate}
\end{enumerate}
\end{proposition}

\begin{proof}[Proof idea]
Use one partition of $[0,+\infty]$ for every vector appearing in the convexity
inequality and push the mass of each cell to a common representative.  On
$[0,1)$ the representatives are chosen from below, so the discrete moment bound
is preserved.  The Shannon branch is kept separate, and in the exceptional
negative case the point $\alpha_*$ is fixed throughout the approximation.  Each
finite approximation satisfies Lemma~\ref{lem:discrete-measure}.

For fixed dimension, $0\leq H_\alpha(\hat{x})\leq\log d$, including the endpoint
values.  Dominated convergence therefore gives pointwise convergence of the
entropy integrals and of their exponentials, and a pointwise limit preserves the
required convexity or concavity inequality.  This is the only convergence input
needed here.  Notice especially that the all-lower branch with $t<0$ requires no
integrability of $\alpha/(1-\alpha)$ near $1$.
\end{proof}

\subsection{Monotonicity under mixing channels}
\label{sec:mixing-channels}

As mentioned at the beginning of this section, monotonicity under mixing channels
acting on the first register $X$ always holds.  Mixing on $X$ increases every
$H_\alpha(P_X)$ and hence the integrated entropy.  Consequently
$\Phi_{t,\tau}$ increases for $t>0$ and decreases for $t<0$; after applying
$(1/t)\log$, the associated conditional entropy increases in both cases.  This
gives the appropriate positive- and negative-homomorphism order directions.
Proposition~A.1 and
Lemma~A.3 of Ref.~\cite{rubboli2026complete} then convert the columnwise
concavity or convexity established above into monotonicity under the conditioning
operation and lift the two components to the original conditionally mixing
channels.  For derivations, the corresponding statements are Lemmas~A.2 and~A.4
of that reference.

\subsection{Sufficient conditions for tropical homomorphisms}

We next take the tropical limits.  Monotonicity is preserved under pointwise
limits, so the temperate conditions immediately provide the admissible endpoint
families.

\begin{proposition}
\label{prop:sufficient-tropical}
Let $\tau$ be a probability measure on $[0,+\infty]$.  The function
\begin{equation}
\label{eq:TropicalForm}
P_{XY}\longmapsto
\min_{\substack{y\in\Y\\P_Y(y)>0}}
\int_{[0,+\infty]}H_\alpha(P_{X|Y=y})\,\d\tau(\alpha)
\end{equation}
is monotone under conditionally mixing channels if either
\begin{enumerate}[label=(\alph*)]
\item $\operatorname{supp}(\tau)\subseteq[0,1]$, with no moment-finiteness
condition, or
\item there is $\alpha_*\in(1,+\infty]$ such that
$\operatorname{supp}(\tau)\subseteq[0,1]\cup\{\alpha_*\}$,
$\tau(\{1\})=0$, both one-sided integrals are finite, and
\begin{equation}
\int_{[0,1)}\frac{\alpha}{1-\alpha}\,\d\tau(\alpha)
+\int_{(1,+\infty]}\frac{\alpha}{1-\alpha}\,\d\tau(\alpha)\leq0.
\end{equation}
\end{enumerate}
\end{proposition}

\begin{proof}[Proof idea]
Use
\begin{align}
\label{eq:limit-tropical}
&\lim_{t\to-\infty}\frac1t\log\!\left(
\sum_{y\in\Y}P_Y(y)
\exp\!\left(t\int H_\alpha(P_{X|Y=y})\,\d\tau(\alpha)\right)
\right)\\
&\hspace{8em}=
\min_{\substack{y\in\Y\\P_Y(y)>0}}
\int H_\alpha(P_{X|Y=y})\,\d\tau(\alpha).
\end{align}
The all-lower case is the limit of Proposition~\ref{prop:sufficient-conditions}(2a).
If the signed moment is strictly negative, the exceptional case holds for all
sufficiently large $|t|$.  At equality, move an arbitrarily small amount of mass
to $\alpha_*$, apply the strict case, and take a final pointwise limit.  This
retains both normalization and the boundary of the parameter range.
\end{proof}

\begin{proposition}
\label{prop:sufficient-tropical-positive}
The function
\begin{equation}
P_{XY}\longmapsto
\max_{\substack{y\in\Y\\P_Y(y)>0}}H_0(P_{X|Y=y})
\end{equation}
is monotone under conditionally mixing channels.
\end{proposition}

\begin{proof}[Proof idea]
This is the $t\to+\infty$ limit of the positive temperate branch with
$\tau$ concentrated at $\alpha=0$:
\begin{equation}
\lim_{t\to+\infty}\frac1t\log\!\left(
\sum_{y\in\Y}P_Y(y)e^{tH_0(P_{X|Y=y})}\right)
=\max_{\substack{y\in\Y\\P_Y(y)>0}}H_0(P_{X|Y=y}).
\end{equation}
\end{proof}

\begin{remark}
Unless $\alpha=0$, the positive temperate moment bound is eventually violated as
$t\to+\infty$.  Thus the positive tropical family contains only the
$\alpha=0$ point mass.  In contrast, every point mass is allowed in the negative
tropical family, subject to the two alternatives in
Proposition~\ref{prop:sufficient-tropical}.
\end{remark}

\subsection{Sufficient conditions for derivations}

We now turn to the derivations.

\begin{proposition}
\label{prop:sufficient-derivations}
The function
\begin{equation}
P_{XY}\longmapsto
\sum_{y\in\Y}P_Y(y)H_\alpha(P_{X|Y=y})
\end{equation}
is monotone under conditionally mixing channels if $\alpha\in[0,1]$.
\end{proposition}

\begin{proof}[Proof idea]
For the point mass at the actual parameter $\alpha$, take $t\to0^-$ in the
all-lower branch of Proposition~\ref{prop:sufficient-conditions}.  The finite
log-sum-exp identity
\begin{equation}
\lim_{t\to0}\frac1t\log\!\left(
\sum_{y\in\Y}P_Y(y)e^{tH_\alpha(P_{X|Y=y})}\right)
=\sum_{y\in\Y}P_Y(y)H_\alpha(P_{X|Y=y})
\end{equation}
passes monotonicity to the limit.  Equivalently, this is the concavity of the
perspective $\vec{x}\mapsto\|\vec{x}\|_1H_\alpha(\hat{x})$ used in Lemma~A.2
of Ref.~\cite{rubboli2026complete}.
\end{proof}

These results identify the admissible homomorphisms and derivations.  With the
normalizations of Definition~6.1 in Ref.~\cite{rubboli2026complete}, they give
the families in Eqs.~(122)--(125) of that reference.

\section{Necessary conditions for the parameters}
\label{sec:necessary-conditions-measures}

In this section, we show that the sufficient conditions on the parameters derived
in the previous section are also necessary.  No assumption is made that the
measure vanishes at a prescribed rate as $\alpha\to1$, and no global finiteness
condition is placed on the weighted integral.  In the positive branch, concavity
forces the lower one-sided moment to be finite and bounded.  In the exceptional
negative branch, convexity forces both one-sided moments to be finite, whereas
the all-lower negative branch requires no such finiteness.

Our approach consists in identifying explicit perturbation directions and
computing the signs of the second derivatives along them.  Outside the ranges of
Section~\ref{sec:sufficient-conditions-measure}, the resulting curvature has the
wrong sign.  Sharp counterexamples require increasingly high dimensions, but the
underlying idea is the finite-dimensional multivariate monomial
\begin{equation}
\label{eq:multivariate-divergences-function}
\vec{v}\longmapsto v_1^{\beta_0}v_2^{\beta_1}\cdots v_d^{\beta_d},
\qquad \beta_0=1-\sum_{i=1}^d\beta_i,
\end{equation}
whose curvature criterion is Lemma~\ref{lem:Matrix-majorization}.  In the
discrete entropy expression, the correspondence is
\begin{equation}
\beta_i\longleftrightarrow
\tau(\{\alpha_i\})t\frac{\alpha_i}{1-\alpha_i},
\qquad
v_i\longleftrightarrow
\left(\sum_jx_j^{\alpha_i}\right)^{1/\alpha_i}.
\end{equation}
The latter quantities are coupled through $\vec{x}$.  Blocks whose cardinalities
and weights scale with a dimension parameter make different blocks dominate on
different ranges of $\alpha$.  For example, appending $d$ entries of weight
$1/d$ gives
\begin{equation}
\left(\sum_jy_j^\alpha\right)^{1/\alpha}
=\left(\sum_{j=1}^k x_j^\alpha+d^{1-\alpha}\right)^{1/\alpha},
\end{equation}
which diverges for $\alpha<1$ and converges to
$(\sum_{j=1}^k x_j^\alpha)^{1/\alpha}$ for $\alpha>1$.

There is one analytic point that is important even in this abbreviated account.
For each localized path used below, its support is fixed near the perturbation
point and one maximizing coordinate remains fixed there.  On these paths the
endpoint-aware kernels and their first two path derivatives extend continuously
through $0$, $1$, and $+\infty$; in particular, the apparent pole at $1$ is
removable.  One differentiates the finite-dimensional kernels before integrating,
then truncates the high-dimensional analysis on the two sides of $1$ and passes
to the endpoint only afterward.  Domination by total variation and monotone
convergence make the required moment bounds conclusions of the argument.  Thus
no derivative is exchanged with an already limiting piecewise envelope, and no
separate finiteness assumption is needed.

\subsection{Necessary conditions; temperate case}

The temperate case corresponds to the bulk entropies $H_{t,\tau}$ in Eq.~(122) of
Ref.~\cite{rubboli2026complete}.  As above, it is sufficient to study
Eq.~\eqref{eq:conditionfunct}.  For notational convenience, combine the weight
$t$ with the probability measure $\tau$ and write
\begin{equation}
\mu:=t\tau.
\end{equation}
Whenever the two one-sided integrals are finite, we also have the coefficient
\begin{equation}
\label{eq:definition-beta0}
\beta_0=1-
\int_{[0,+\infty]\setminus\{1\}}
\frac{\alpha}{1-\alpha}\,\d\mu(\alpha),
\end{equation}
where the integral is the sum of its parts on $[0,1)$ and $(1,+\infty]$.

Whenever a signed moment is written as a single ordinary real number below, it is
the sum of the two one-sided integrals after their finiteness has been established.
Finiteness is not a blanket hypothesis on $\mu$: before it is known, the argument
works with endpoint-separated truncations.  The localized estimates show that if
convexity leaves any support above $1$, the relevant one-sided moments must be
finite.  If all support lies in $[0,1]$, no such condition is imposed.

We first treat $\mu>0$, and then $\mu<0$.  In the negative case the Shannon atom,
the signed moment, and the number of upper support points are isolated in separate
steps, exactly matching the alternatives in
Proposition~\ref{prop:sufficient-conditions}.

\begin{proposition}
\label{prop:first-measure}
Let $\mu$ be a positive finite measure on $[0,+\infty]$.  The function
\begin{equation}
\vec{x}\longmapsto\|\vec{x}\|_1
\exp\!\left(\int_{[0,+\infty]}H_\alpha(\hat{x})\,
\d\mu(\alpha)\right)
\end{equation}
is concave only if
$\operatorname{supp}(\mu)\subseteq[0,1]$, $\mu(\{1\})=0$, and
\begin{equation}
\int_{[0,1)}\frac{\alpha}{1-\alpha}\,\d\mu(\alpha)\leq1.
\end{equation}
\end{proposition}

\begin{proof}[Proof idea]
For $p\in(0,1)$, use the two-block point and traceless direction
\begin{align}
\vec{x}&=\left[
\underbrace{\frac p d,\ldots,\frac p d}_{d\ \mathrm{times}},
\underbrace{\frac{1-p}{d^2},\ldots,\frac{1-p}{d^2}}_{d^2\ \mathrm{times}}
\right],\\
\vec{v}&=\left[
\underbrace{-\frac1d,\ldots,-\frac1d}_{d\ \mathrm{times}},
\underbrace{\frac1{d^2},\ldots,\frac1{d^2}}_{d^2\ \mathrm{times}}
\right].
\end{align}
As $d\to\infty$, the two blocks localize the ranges below and above $1$.  Any
Shannon mass produces a strictly positive leading curvature.  Upper mass gives a
coefficient $\beta<0$, hence $\beta^2-\beta>0$, and a lower moment exceeding
$1$ gives the same wrong sign with $\beta>1$.  Each case contradicts concavity.
The truncation procedure described above justifies the limits and simultaneously
shows that the lower moment is finite and bounded by $1$.
\end{proof}

Next, let $\mu$ be negative.

\begin{proposition}
\label{prop:no-alpha-one}
Let $\mu$ be a negative finite measure with $\mu(\{1\})\neq0$.  Then the function
\begin{equation}
\vec{x}\longmapsto\|\vec{x}\|_1
\exp\!\left(\int_{[0,+\infty]}H_\alpha(\hat{x})\,
\d\mu(\alpha)\right)
\end{equation}
is convex only if $\operatorname{supp}(\mu)\subseteq[0,1]$.
\end{proposition}

\begin{proof}[Proof idea]
Assume that both the Shannon atom and upper support are present.  If
$\mu(\{1\})<0$ and $\beta_1>0$ is a localized upper contribution, the
two-block direction is scaled as $(-\beta_1/\mu(\{1\}),1)$.  The first
logarithmic derivative then tends to zero and the second tends to
$-\beta_1<0$.  This is a
genuinely concave direction and contradicts convexity.  Working first on
truncated ranges makes the calculation finite; the endpoint limits are taken only
after the curvature sign is fixed.
\end{proof}

\begin{proposition}
\label{prop:second-measure}
Let $\mu$ be a negative finite measure with $\mu(\{1\})=0$, and suppose that the
function
\begin{equation}
\vec{x}\longmapsto\|\vec{x}\|_1
\exp\!\left(\int_{[0,+\infty]}H_\alpha(\hat{x})\,
\d\mu(\alpha)\right)
\end{equation}
is convex.  If $\operatorname{supp}(\mu)$ meets $(1,+\infty]$, then both
one-sided integrals are finite and
\begin{equation}
\label{eq:negative-moment-necessary}
\int_{[0,1)}\frac{\alpha}{1-\alpha}\,\d\mu(\alpha)
+\int_{(1,+\infty]}\frac{\alpha}{1-\alpha}\,\d\mu(\alpha)\geq1.
\end{equation}
Consequently, a finite signed sum smaller than $1$ is compatible with convexity
only if $\operatorname{supp}(\mu)\subseteq[0,1]$.
\end{proposition}

\begin{proof}[Proof idea]
First use separated upper neighborhoods with stationary block velocities.  If
two distinct upper portions survived, their limiting second derivative would be
negative, so convexity leaves at most one upper point; its upper moment is
therefore finite.  A three-block construction then controls normalization and
localizes the lower and upper ranges.  The perturbation
is chosen so that the first derivative is exactly zero.  On every truncation away
from $1$, a signed moment smaller than $1$ makes the limiting second derivative
negative.  Convexity rules this out.  Monotone passage through the truncations
forces finiteness of the lower moment and gives
Eq.~\eqref{eq:negative-moment-necessary}.  This is how endpoint localization
replaces the old global finiteness hypothesis.
\end{proof}

\begin{remark}
The points used in these counterexamples need not initially be normalized.  The
function is positively homogeneous on the whole semiring domain, so rescaling by
a positive constant does not change the curvature sign.  Equivalently, the
counterexample can be embedded into a normalized joint probability matrix and
then interpreted in the semiring of Section~3B and lifted using
Proposition~A.1 and Lemma~A.3 of Ref.~\cite{rubboli2026complete}.
\end{remark}

\begin{remark}
A nontraceless direction can likewise be embedded into a trace-preserving
two-column perturbation.  If the first column changes as
$\vec{x}+\lambda\vec{v}$, change the second as
\begin{equation}
\vec{x}'-\lambda\hat{x}'\sum_i v_i.
\end{equation}
The opposite sign is essential: it keeps the total weight fixed.  Additivity over
the conditioning register makes the second column affine in $\lambda$, so it
does not alter the curvature obstruction.
\end{remark}

Finally, the upper support cannot contain two distinct points.

\begin{proposition}
\label{prop:third-measure}
Let $\mu$ be a negative finite measure with $\mu(\{1\})=0$, and suppose that the
function
\begin{equation}
\vec{x}\longmapsto\|\vec{x}\|_1
\exp\!\left(\int_{[0,+\infty]}H_\alpha(\hat{x})\,
\d\mu(\alpha)\right)
\end{equation}
is convex.  Then either $\operatorname{supp}(\mu)\subseteq[0,1]$, or there is a
unique $\alpha_*\in(1,+\infty]$ such that
\begin{equation}
\operatorname{supp}(\mu)\subseteq[0,1]\cup\{\alpha_*\}.
\end{equation}
In the second case the two one-sided moments are finite and satisfy
Eq.~\eqref{eq:negative-moment-necessary}.
\end{proposition}

\begin{proof}[Proof idea]
If two disjoint portions of the support lie above $1$, localize them in two
different dominant blocks.  Their positive localized weights may be denoted
$\beta_1,\beta_2>0$.  The block velocities are chosen to cancel the first
derivative, while the limiting second derivative is
\begin{equation}
-\frac1{\beta_1}-\frac1{\beta_2}<0.
\end{equation}
This contradicts convexity.  Shrinking the two localization windows shows that
the upper support has at most one point.  The finiteness and signed inequality in
the exceptional branch are the conclusion of the joint truncation argument in
Proposition~\ref{prop:second-measure}.
\end{proof}

Thus the last two propositions cover all negative measures without a global
finiteness assumption; the all-lower branch still carries no moment condition.

The same stationary localization proves necessity at the negative tropical
endpoint: convexity is replaced by quasiconvexity, so one first imposes exact
stationarity and then uses the negative second derivative.  The result is exactly
the two alternatives of Proposition~\ref{prop:sufficient-tropical}, with the
right-hand side of the signed moment bound equal to $0$.  At the positive
tropical endpoint, strong max-quasiconcavity on simplex faces forces
$\tau$ to be concentrated at $\alpha=0$: strong max-quasiconcavity would force
the integrated entropy to be constant on each fixed-support face, whereas
comparing a nonuniform point with the uniform point gives a strict $H_\alpha$ gap
for every $\alpha>0$.
Thus the negative tropical endpoint has exactly the measures of
Proposition~\ref{prop:sufficient-tropical}, whereas the positive tropical endpoint
is only $H_{+\infty,0}$ in Eq.~(125) of Ref.~\cite{rubboli2026complete}.

\subsection{Necessary conditions for derivations}
\label{sec:necessary-derivations}

Next, we turn to the necessary conditions for the extremal derivations.  According
to Eq.~(91) and Proposition~3.7 of Ref.~\cite{rubboli2026complete}, these are
\begin{equation}
H_{0,\alpha}(P_{XY})=
\sum_{y\in\Y}P_Y(y)H_\alpha(P_{X|Y=y}).
\end{equation}
It is therefore sufficient to consider the concavity of
\begin{equation}
\vec{x}\longmapsto\|\vec{x}\|_1H_\alpha(\hat{x}),
\end{equation}
as required by Lemma~A.2 of Ref.~\cite{rubboli2026complete}.  We now show that
every $\alpha>1$, including $+\infty$, is excluded.

\begin{proposition}
\label{prop:measure-derivation}
The function
\begin{equation}
\vec{x}\longmapsto\|\vec{x}\|_1H_\alpha(\hat{x})
\end{equation}
is concave only if $\alpha\in[0,1]$.
\end{proposition}

\begin{proof}[Proof idea]
Use the two-block point and direction from Proposition~\ref{prop:first-measure}.
For $\alpha>1$, put
\begin{equation}
\beta=\frac{\alpha}{1-\alpha}<0.
\end{equation}
After the large-dimensional localization, the second derivative at the origin
tends, up to a positive factor, to
\begin{equation}
-\frac{\beta}{p^2}>0.
\end{equation}
This is a convex direction and contradicts the required concavity.  The endpoint
$\alpha=+\infty$ follows from the same localization by taking the endpoint limit.
\end{proof}

Importing the upstream representation of derivations from Proposition~3.7 and
Eq.~(91) of Ref.~\cite{rubboli2026complete}, and combining it with
Proposition~\ref{prop:sufficient-derivations} and the semiring lift, gives the
final statement.

\begin{proposition}
\label{prop:Derivations}
The set of derivations of the conditional majorization semiring $S$ at
$\|\cdot\|$ is spanned by the quantities $H_{0,\alpha}$ of Definition~6.1 in
Ref.~\cite{rubboli2026complete}, with $\alpha\in[0,1]$.
\end{proposition}

The span assertion in Proposition~\ref{prop:Derivations} is imported from that
upstream representation.  The sufficiency and curvature arguments in this note
supply the exact restriction $\alpha\in[0,1]$ for its extremal candidates.

Together, the temperate, tropical, and derivation statements give the exact
parameter ranges for the concrete candidate families in Eqs.~(122)--(125) of
Ref.~\cite{rubboli2026complete}; see also Definitions~6.2 and~6.3 there for the
parameter sets.  They supply the parameter-classification input used in
Theorem~7.1 of that paper.  The upstream representation theorem and the later
operational applications remain results of the cited paper rather than claims of
this standalone note.

\bibliographystyle{ultimate}
\bibliography{bibliography}

\end{document}
